\documentclass[conference]{IEEEtran}
\usepackage{graphicx}
\usepackage{cite}
\usepackage{microtype}
\usepackage{amsmath,amssymb}
\usepackage{booktabs}
\begin{document}
\title{CIFAR-10 Classification: Augmentation and Optimization Experiments}
\author{Must\\Software Engineering Student\\Institution}
\maketitle
\begin{abstract}
We design and evaluate a convolutional neural network for CIFAR-10 classification exploring two paths: data augmentation (CutMix/Cutout) and optimization strategies (SGD vs Adam with LR scheduling). We report experimental comparisons, per-class performance, and visual explainability results.
\end{abstract}

\section{Introduction}
This project compares augmentation and optimization choices for CIFAR-10. Our goal is to justify design choices and analyze their impact on generalization and robustness.

\section{Model Architecture}
We adopt a ResNet-18 backbone adjusted for CIFAR-10 (3x32 input): first convolution swapped to 3x3, no maxpool, final FC adapted to 10 classes. Model has ~11M parameters.

\section{Experiments \& Results}
Two experiments were conducted:
\begin{itemize}
  \item \textbf{Augmentation:} baseline (random crop+flip) vs CutMix (mixing images). CutMix improved validation accuracy and reduced overfitting observed in training curves.
  \item \textbf{Optimization:} SGD with momentum + CosineAnnealing LR outperformed Adam for final validation accuracy in our runs; Adam converged faster but to a slightly lower peak.
\end{itemize}

We report accuracy, confusion matrix, and per-class accuracy in the Colab notebook. Grad-CAM visualizations highlight class-discriminative regions, supporting model interpretability.

\section{Conclusion}
Augmentation (CutMix) and careful optimization (SGD+LR scheduling) improved CIFAR-10 performance. Detailed plots and code are provided in the Colab notebook.

\section*{References}
[1] H. Zhang et al., ``mixup: Beyond empirical risk minimization'', ICLR 2018.\\
[2] Y. Yun et al., ``CutMix: Regularization Strategy to Train Strong Classifiers'', ICCV 2019.
\end{document}
